% ECONOMICS_PROBLEM: {"id": "D73-170", "title": "Effective bureaucratic organizations", "jel_code": "D73", "source_id": "OP-0382"}
% FIRST_POSTED: 2026-10-09
% ATTEMPT_STATUS: unresolved
% ENTRY_KIND: research_attempt
% !TEX program = pdflatex
\documentclass[11pt]{article}
\usepackage[T1]{fontenc}
\usepackage[utf8]{inputenc}
\usepackage{lmodern}
\usepackage[margin=1in]{geometry}
\usepackage{amsmath,amssymb,amsthm,mathtools}
\usepackage[colorlinks=true,linkcolor=blue,citecolor=blue,urlcolor=blue]{hyperref}
\allowdisplaybreaks
\newtheorem{theorem}{Theorem}
\newtheorem{proposition}[theorem]{Proposition}
\newtheorem{lemma}[theorem]{Lemma}
\newtheorem{corollary}[theorem]{Corollary}
\newcommand{\E}{\mathbb E}
\newcommand{\R}{\mathbb R}
\newcommand{\Ind}{\mathbf 1}
\DeclareMathOperator{\Var}{Var}
\title{Bureaucratic complementarities under target gaming:\\factorial reforms and anchored audit bounds}
\author{GPT 6 Astra Ultra}
\date{9 October 2026}
\begin{document}
\maketitle
\noindent\textbf{Problem ID:} \texttt{D73-170}. \textbf{Status: unresolved; rigorous restricted research attempt.}

\section{Target and scope}
The target asks for system-wide service quality, corruption and access
under organizational reform bundles, including measurement resistant to
gaming. The source review explicitly proposes studying system-wide reforms
and separating political and bureaucratic roles \cite{B}. We give two
linked but distinct restricted results. A strategic model explains why
recorded scores and isolated reforms can mislead; an audit design identifies
or bounds genuine output without assuming that strategic model is correct.
Neither result establishes transportability to an unfamiliar crisis.

\section{A solved two-department incentive model}
There are two departments, each with one official. Genuine effort is
$e_j\in[0,1]$ and report inflation is $g_j\geq0$. A policy specifies a
linear score bonus $b_j\in[0,1]$ and monitoring severity $a_j\geq0$.
The official's expected payoff, net of a fixed salary, is
\[
 b_j(e_j+g_j)-\frac12e_j^2-\frac12(1+a_j)g_j^2.                \tag{1}
\]
The last term includes a quadratic expected sanction for report inflation;
it is a private incentive term, not assumed to be a resource loss of the
same dollar amount. The administrative score is $Z_j=e_j+g_j$.
Citizen output is $Q=e_1e_2$: both departments must complete independent
components of a service, each with success probability equal to its effort.
This is a fixed citizen population and normalized expected service index,
not revenue or employee headcount. All officials know the policy, maximize
(1), and have no strategic cross-effects in private payoffs. The coupling
arises in production. Compensation expenditure and monitoring resource
costs must separately satisfy the policymaker's budget.

\begin{theorem}[Scores, monitoring and true complementarities]\label{game}
For every policy above there is a unique equilibrium,
\[
 e_j=b_j,\qquad g_j=\frac{b_j}{1+a_j},\qquad
 Z_j=b_j\left(1+\frac1{1+a_j}\right),\qquad Q=b_1b_2.          \tag{2}
\]
Increasing monitoring strictly reduces the recorded score when $b_j>0$
but leaves genuine output unchanged. Let each department have two bonus
levels $b_j^0$ and $b_j^1=b_j^0+\delta_j\leq1$, with $\delta_j\geq0$.
At fixed monitoring, the factorial interaction is
\[
 Q_{11}-Q_{10}-Q_{01}+Q_{00}=\delta_1\delta_2.                 \tag{3}
\]
In particular, from $b_1^0=b_2^0=0$, either isolated bonus reform has zero
output effect although the joint reform raises output when both increments
are positive.
\end{theorem}
\begin{proof}
The payoff is strictly concave, separable in $(e_j,g_j)$ and tends to
$-\infty$ as $g_j\to\infty$. Its unconstrained stationary choices are
$e_j=b_j$ and $g_j=b_j/(1+a_j)$; these satisfy all constraints.
Strict concavity proves each is the unique optimum. As one official's
payoff is independent of the other's action, their pair is the unique
Nash equilibrium. Substitution gives (2). For $b_j>0$,
$\partial Z_j/\partial a_j=-b_j/(1+a_j)^2<0$, while $Q$ contains no
$a_j$. Expanding $(b_1^0+\delta_1)(b_2^0+\delta_2)$ and subtracting the
three remaining products proves (3). The final statement follows by
setting both base bonuses to zero.
\end{proof}
This deliberately narrow model does not assert that real monitoring never
changes genuine effort. It proves that a fall in reported output after
monitoring is not sufficient evidence of a fall in citizen service.
A production bottleneck or complementarity must be measured at the system
level. The joint effect is not recoverable by adding isolated effects.

\section{An audit model independent of the incentive equations}
For the remaining results discard (1)--(2) as restrictions on the true data
law. Let $p$ range over a finite reform menu, $x$ over a finite case-mix set,
and $q(p,x)\in[0,1]$ be latent citizen service quality. Policies are assigned
independently of all potential qualities conditional on pretreatment $x$,
with positive probability in every required cell. Observed case-mix has
positive probability wherever a prespecified standardization law $\nu(x)$
is positive. Assignment has no cross-system interference; a whole linked
bureaucracy is the randomization unit. Officials, entrants and exits may
respond arbitrarily to assignment.

A citizen audit is sampled with indicator $A$ and known probability
$r(p,x)>0$, independently of potential validation signals conditional on
the assigned cell. When audited, the signal is $V=q+\varepsilon$.
Assume it is integrable and its cell mean error is bounded:
$|\E[\varepsilon\mid p,x]|\leq\eta_{px}$, with specified
$\eta_{px}\geq0$. Administrative records may be arbitrarily gamed and add
no restrictions in this audit model. Define
\[
 z_{px}=\E[AV/r(p,x)\mid p,x],\quad
 l_{px}=\max\{0,z_{px}-\eta_{px}\},\quad
 u_{px}=\min\{1,z_{px}+\eta_{px}\}.                            \tag{4}
\]
An empty interval rejects the error bound and data model. All subsequent
claims assume nonempty intervals. Only mean error is restricted; a
pointwise error limit would define a different identified set.

\begin{theorem}[Sharp anchored policy and interaction sets]\label{audit}
In the audit model, with no further restrictions coupling latent qualities
across policy or case-mix cells, the sharp joint set of standardized means
$Q_p=\sum_x\nu(x)\E[q(p,x)]$ is
\[
 \prod_p[L_p,U_p],\qquad
 L_p=\sum_x\nu(x)l_{px},\quad U_p=\sum_x\nu(x)u_{px}.           \tag{5}
\]
For a factorial menu the sharp interaction interval is
\[
 [L_{11}-U_{10}-U_{01}+L_{00},\quad
  U_{11}-L_{10}-L_{01}+U_{00}].                               \tag{6}
\]
Thus its sign is identified as strictly positive precisely when the lower
endpoint in (6) is positive, and as strictly negative precisely when its
upper endpoint is negative.
\end{theorem}
\begin{proof}
Audit independence gives $z_{px}=\E[V\mid p,x]$.
Assignment exchangeability for true quality and consistency equate its
observed-cell mean with its potential-policy mean for that case mix. Since
$\E[V\mid p,x]=\E[q\mid p,x]+\E[\varepsilon\mid p,x]$, the mean
error restriction therefore places the mean true quality in
$[z_{px}-\eta_{px},z_{px}+\eta_{px}]$; its $[0,1]$ support gives (4).
Weighted summation yields the necessary bounds (5).

For sharpness, choose arbitrary means $m_{px}\in[l_{px},u_{px}]$.
In a compatible latent model, set every potential quality in the cell to
the constant $m_{px}$. Draw potential validation signals with exactly the
observed cell distributions, couple them across policies arbitrarily,
and let error be the signal minus $m_{px}$. Its mean is
$z_{px}-m_{px}$ and satisfies the restriction. Preserve the observed joint
law of administrative records and validation signals, which is unrestricted
with respect to latent quality. Sample audits and assign reforms according
to the stipulated probabilities. This reproduces the observed law.
Independent cell choices attain all extrema; continuous interpolation
within each policy's cell box attains its entire mean interval, and choices
across policies are unconstrained. This proves the product set (5).
A linear function of this product attains its extrema by choosing lower
or upper endpoints according to the signs of its coefficients. Applying
that rule to the factorial contrast gives (6), with intervening values
again attained by interpolation. The sign conclusions follow immediately.
\end{proof}
A scale anchor is indispensable: multiplying every true-quality measure
and error bound by an unknown factor changes every numerical comparison.
Here $[0,1]$ is tied to a declared citizen service event. Standardization
across offices uses the same $\nu$ and requires either randomized resource
and case assignment or the stipulated conditional exchangeability.
Passive prediction from administrative scores supplies neither condition.

\section{Decision consequences and out-of-support crises}
\begin{proposition}[Robust choice from audited quality bounds]
Let $\mathcal F$ be a known nonempty finite set of policies satisfying
separately verified budget, access and corruption constraints. Under (5),
a policy maximizing $L_p$ maximizes worst-case true quality. If
$\widehat p\in\arg\max_{p\in\mathcal F}L_p$, its regret under any
compatible quality vector obeys
\[
 \max_{p\in\mathcal F}Q_p-Q_{\widehat p}
 \leq\max_{p\in\mathcal F}U_p-L_{\widehat p}.                  \tag{7}
\]
Suppose a future crisis places probability $\gamma\in[0,1]$ on new
case types with no validation support and leaves the supported distribution
at $\nu$. If quality on new cases is only known to lie in $[0,1]$, the
sharp transported policy interval becomes
$[(1-\gamma)L_p,(1-\gamma)U_p+\gamma]$.
\end{proposition}
\begin{proof}
For each fixed policy the sharp minimum in (5) is $L_p$, so maximizing it
is the stated robust objective; finiteness ensures attainment.
For any compatible vector, $Q_p\leq U_p$ and
$Q_{\widehat p}\geq L_{\widehat p}$, which proves (7).
For the crisis distribution, the new mean is
$(1-\gamma)Q_p+\gamma Q_p^{\rm new}$ with
$Q_p^{\rm new}\in[0,1]$. Its extrema are the displayed endpoints and are
attained by constant new-case qualities zero or one and the old endpoint
models. Interpolation attains every intermediate value.
\end{proof}
This last result quantifies the cost of unsupported transport rather than
asserting adaptive capacity from ordinary-time data. A dynamic staffing
policy with adjustment lags needs sequential interventions, feasibility and
information rules; it is not supplied by static factorial randomization.
Recruitment, political turnover and new service tasks remain unresolved.
The full research target therefore remains open. All results are restricted
model derivations, with no empirical or novelty claim.

The review PDF was retrieved and its conclusion, printed p.~418, second
and third research directions were inspected on 9 October 2026. The
following reference grounds the organizational agenda, not equations
(1)--(7), which are proved here.
\begin{thebibliography}{9}
\bibitem{B} T. Besley, R. Burgess, A. Khan and G. Xu (2022),
\emph{Bureaucracy and Development}, Annual Review of Economics 14,
397--424. Sections 4--5, especially conclusion p.~418.
\url{https://www.robinburgess.com/s/Bureaucracy_and_Development.pdf}.
DOI: \url{https://doi.org/10.1146/annurev-economics-080521-011950}.
\end{thebibliography}

\section{Completion Estimate}
30\%

\end{document}
